\section{Tohoku Case Study and Data}

The event source is the 2011 off-Pacific-coast Tohoku earthquake. Repository preprocessing uses approved finite-fault input, with USGS and published source reconstructions retained as authority records \autocite{usgs2011tohokuEvent,usgsFiniteFaults,hayes2011rapid,fujii2011tsunamisource}. Source sensitivity remains an early calibration target because Tohoku-specific studies show sensitivity to source geometry, kinematics and dispersive/model-form assumptions \autocite{saito2014dispersion,baba2017farfield,wei2013japan}.

Terrain and bathymetry are treated as measured or fidelity/model-form inputs rather than free calibration variables. The G6 baseline uses ETOPO 2022 as the accepted available terrain source, while higher-resolution J-EGG500, GEBCO and GSI products remain documented candidates for later verification and validation work \autocite{noaa2022etopo,gebco2026grid,jodcJegg500,gsiDemDownload}.

Observation sources are not yet approved for calibration execution. The project expects offshore waveforms, arrival times, run-up heights and inundation extents to be ingested under an observation-approval process. Mori et al. provide a principal post-event survey authority, but the survey type, spatial density and tide-correction provenance must be respected \autocite{mori2012nationwide}.
